\section{本章小结}
\label{sec:qn-summary}

本章开头的问题是：Hessian算不起、存不下时，能否从迭代本身得到曲率信息。答案是相邻两点的梯度差。式~\eqref{eq:qn-average-hessian}说明它是平均Hessian作用在位移上的结果，由此得到割线条件。这个条件在$n=1$时确定了曲率，$n>1$时还缺$n(n-1)/2$个方程，所以需要补充对称、低秩和曲率的要求。

不同的补充要求产生了不同的方法。DFP和BFGS各用两个秩一项，分别在矩阵和逆矩阵上度量改动。两者都在曲率条件$\langle\bm{s},\bm{y}\rangle>0$下保持正定，Wolfe线搜索使这个条件对任意光滑目标成立，跳过和阻尼是线搜索之外的备用办法。

SR1只用一个秩一项，由割线条件唯一确定。它在二次目标上有更强的遗传性质，能记录负曲率，但分母可能近零，矩阵也可能不定，所以需要跳过规则，且方向分析不能照搬，要交给有界的模型处理。

L-BFGS用双循环递推在不保存矩阵的情况下计算$\bm{W}^{(t)}\bm{g}^{(t)}$，把存储从$n^2/2$降到$2rn$。记忆长度决定了方向的质量和每步的计算量，初始缩放决定了单位步长是否合适。

\begin{BookInlineTable}
  \centering
  \begin{tabular}{llll}
    \toprule
    方法 & 新增信息 & 保持的性质 & 单步主要开销\\
    \midrule
    BFGS & 秩二修正 & 曲率条件下正定 & $O(n^2)$\\
    DFP & 秩二修正 & 曲率条件下正定 & $O(n^2)$\\
    SR1 & 秩一修正 & 可表示负曲率 & $O(n^2)$\\
    L-BFGS & 最近$r$对 & 曲率条件下正定 & $O(rn)$\\
    \bottomrule
  \end{tabular}
  \caption{本章四种拟牛顿方法的比较。}
  \label{tab:qn-summary}
\end{BookInlineTable}

收敛结论要分开看。BFGS在强凸目标和Wolfe线搜索下全局收敛，在更强的正则性下超线性收敛，后者由Dennis--Mor\'e条件刻画，它只要求沿步方向准确。L-BFGS在强凸目标上线性收敛，完整记忆的超线性结论不能自动移给它。一般非凸目标上，这些保证都不成立，SR1在一般目标上的局部结论也有额外条件。

拟牛顿矩阵的好坏，最终还要看模型在多大范围内可信。线搜索只沿一个方向试探，不定模型又没有极小点；第~\ref{chap:trust-region}章研究的信赖域，把候选步和接受判断放在一起处理，并且既可以用精确Hessian，也可以用本章的拟牛顿矩阵。
